\documentclass[10pt,twocolumn]{article}
\usepackage[a4paper,margin=1.8cm]{geometry}
\usepackage{kotex}
\usepackage{amsmath,amssymb}
\usepackage{graphicx}
\usepackage{booktabs}
\usepackage{hyperref}
\usepackage{enumitem}
\usepackage{titlesec}
\setlength{\columnsep}{0.7cm}
\setlength{\parskip}{0.3em}
\setlength{\parindent}{1em}
\titleformat{\section}{\bfseries\normalsize}{\thesection.}{0.4em}{}
\titleformat{\subsection}{\bfseries\normalsize}{\thesubsection}{0.4em}{}
\hypersetup{colorlinks=true,linkcolor=black,citecolor=black,urlcolor=blue}

\title{\textbf{비침습 입력 기반 대사증후군 위험 분류:\\생활패턴 스무딩 피처, 라벨 전이, 임계값 최적화 기반 XGBoost 연구}}
\author{CapStone A Team}
\date{}

\begin{document}
\maketitle

\begin{abstract}
본 연구는 혈액검사 없이 수집 가능한 비침습 입력(연령, 성별, 키, 체중, 허리둘레, 영양섭취 상태)으로 대사증후군 위험을 분류하는 머신러닝 파이프라인을 제안한다.
데이터셋 간 직접 병합이 어려운 상황을 해결하기 위해 라벨 전이(label transfer) 전략과 생활패턴 스무딩 피처를 도입하였다.
XGBoost를 GridSearchCV(ROC-AUC 기준)로 최적화한 결과, 기본 임계값(0.50)에서 Accuracy 0.8109, F1 0.5404, ROC-AUC 0.8664를 보였고,
F1 최적 임계값(0.36)에서 F1 0.6530, Recall 0.7571로 개선되었다.
Recall 우선 정책(0.23, precision$\ge$0.50)에서는 Recall 0.8603을 달성했다.
본 결과는 선별 서비스에서 모델 선택뿐 아니라 임계값 운영정책이 핵심임을 보여준다.
\end{abstract}

\noindent\textbf{주요어:} 대사증후군, 비침습 분류, 라벨 전이, 생활패턴 스무딩, XGBoost, 임계값 최적화

\section{서론}
대사증후군은 심혈관계 질환 및 제2형 당뇨병의 선행 위험군 식별에 중요한 임상 개념이다.
그러나 실제 사용자 환경에서는 혈액검사 기반 정밀 진단 이전 단계에서, 비침습 정보 기반 선별 도구가 먼저 요구된다.
본 연구는 이질 코호트 환경에서 비침습 입력만으로 대사증후군 위험을 분류하는 실무형 파이프라인을 제안한다.

본 연구의 핵심 질문은 다음과 같다.
\begin{enumerate}[leftmargin=1.2em]
\item 직접 병합이 어려운 건강/영양 데이터 간 정보 연결이 가능한가?
\item 생활패턴 스무딩 피처가 분류 성능을 개선하는가?
\item 선별 목적에서 고정 임계값(0.50)보다 정책형 임계값이 유리한가?
\end{enumerate}

\section{관련 연구}
대사증후군 기준은 국제 가이드라인 및 국내 적용 기준을 바탕으로 사용하였다.
비침습 선별의 장점은 접근성과 반복성이나, 입력 정보량 제약으로 인해 피처 공학과 threshold 설계가 성능을 좌우한다.
본 연구는 라벨 전이, 스무딩 피처, 운영 임계값 정책을 통합적으로 다룬다는 점에서 차별성을 가진다.

\section{데이터 및 방법}
\subsection{데이터}
건강 데이터는 \texttt{NHANES\_2017\_2023.csv}, 영양 데이터는 \texttt{nutrient\_2019.csv}를 사용하였다.
학습/평가 분할은 stratified 8:2이며, train 38174건, test 9544건이다.

\subsection{라벨 전이와 스무딩 피처}
\texttt{nutrient\_2019}의 다중 식사기록을 개인 단위로 집계해 분포 통계(평균/표준편차/IQR), 과잉/부족/균형 일수 비율, 영양밀도 피처를 산출했다.
이후 KNN 기반 라벨 전이로 NHANES 대상자에 영양 상태 라벨과 확률(\texttt{nutrition\_prob\_2})을 부여했다.
핵심 상호작용 피처는 다음과 같다.
\[
\texttt{waist\_x\_excess\_prob}=\texttt{HE\_wc}\times \texttt{nutrition\_prob\_2}
\]

\subsection{타깃 정의}
다음 5개 기준 중 3개 이상 충족 시 대사증후군 양성(1)으로 라벨링했다.
\begin{itemize}[leftmargin=1.2em]
\item 복부비만: 남 $\ge$90cm, 여 $\ge$80cm
\item 중성지방: TG $\ge$150
\item HDL 저하: 남 $<$40, 여 $<$50
\item 혈압: SBP $\ge$130 또는 DBP $\ge$85
\item 공복혈당: glucose $\ge$100
\end{itemize}

\subsection{모델 및 평가}
XGBoost를 사용하고, 결측치는 KNNImputer로 보정했다.
GridSearchCV(5-fold, ROC-AUC) 최적 파라미터는
\texttt{colsample\_bytree=0.8}, \texttt{learning\_rate=0.01},
\texttt{max\_depth=5}, \texttt{n\_estimators=200}, \texttt{subsample=0.8}이다.
임계값은 0.10--0.90 범위를 0.01 간격으로 탐색했다.

\section{결과}
\subsection{정책별 성능}
\begin{table}[h]
\centering
\caption{임계값 정책별 성능(Full 모델)}
\begin{tabular}{lccccc}
\toprule
Policy & Thr. & Acc. & Prec. & Rec. & F1 \\
\midrule
Default & 0.50 & 0.8109 & 0.6319 & 0.4720 & 0.5404 \\
F1-opt & 0.36 & 0.8105 & 0.5740 & 0.7571 & 0.6530 \\
Recall-opt & 0.23 & 0.7687 & 0.5052 & 0.8603 & 0.6366 \\
\bottomrule
\end{tabular}
\end{table}

기본 임계값은 정밀도 측면에서 유리하지만 재현율이 낮다.
F1 최적 정책은 정밀도 감소를 제한하면서 재현율을 크게 개선한다.
Recall 우선 정책은 미탐 최소화 시나리오에서 효과적이다.

\subsection{A/B 비교}
\begin{table}[h]
\centering
\caption{Full vs Light 비교}
\begin{tabular}{lcccc}
\toprule
Model & \#Feat. & Default F1 & AUC & Best F1 \\
\midrule
Full & 18 & 0.5404 & 0.8664 & 0.6530 \\
Light & 11 & 0.5364 & 0.8662 & 0.6517 \\
\bottomrule
\end{tabular}
\end{table}

Full 모델이 F1과 AUC에서 근소 우세하여 최종 운영 모델로 선정하였다.

\subsection{연령대별 분석}
연령대 분석에서는 고연령군으로 갈수록 Recall/F1이 상승하는 경향이 관찰되었다.
이는 단일 임계값 운영의 한계를 시사하며, 향후 연령군별 threshold 조정 가능성을 제안한다.

\section{논의}
본 연구는 이질 코호트 간 직접 병합이 어려운 환경에서 라벨 전이 기반 접근이 실무적으로 유효함을 보였다.
또한 동일 모델에서도 임계값 정책에 따라 오류 구조가 크게 달라져, 운영 단계의 의사결정 규칙이 모델 구성의 일부가 되어야 함을 확인했다.

생활패턴 스무딩 피처의 개선폭은 소폭이지만 일관성이 있었다.
이는 고잡음/부분관측 환경에서 과도한 복잡도보다 안정적 분포정보가 유효할 수 있음을 의미한다.

\section{한계 및 윤리}
주요 한계는 다음과 같다.
\begin{enumerate}[leftmargin=1.2em]
\item 라벨 전이 과정에서 전이 오차가 누적될 수 있음
\item 단일 홀드아웃 기반 threshold 최적화의 재현성 한계
\item 짧은 기록일수 분포로 장기 식습관 반영이 제한됨
\end{enumerate}
본 모델은 연구/건강관리 참고용이며 의료 진단, 치료, 처방을 대체하지 않는다.

\section{결론}
비침습 입력 기반 대사증후군 선별에서 라벨 전이와 생활패턴 스무딩은 유의미한 성능 확보에 기여했다.
운영 기준으로는 Full 모델의 F1 최적 임계값(0.36)을 기본 정책으로, 미탐 최소화가 필요한 상황에서는 Recall 우선 임계값(0.23)을 보조 정책으로 사용할 수 있다.

\section*{References}
\begin{enumerate}[leftmargin=1.2em]
\item Expert Panel on Detection, Evaluation, and Treatment of High Blood Cholesterol in Adults. Executive summary of the third report of the NCEP Adult Treatment Panel III. \textit{JAMA}. 2001;285(19):2486--2497.
\item Alberti KGMM, Eckel RH, Grundy SM, et al. Harmonizing the metabolic syndrome: a joint interim statement of major international societies. \textit{Circulation}. 2009;120(16):1640--1645.
\item Ministry of Health and Welfare (Korea), KDCA. KNHANES documentation.
\item Pedregosa F, Varoquaux G, Gramfort A, et al. Scikit-learn: Machine Learning in Python. \textit{JMLR}. 2011;12:2825--2830.
\item Chen T, Guestrin C. XGBoost: A scalable tree boosting system. \textit{KDD}. 2016:785--794.
\item Troyanskaya O, Cantor M, Sherlock G, et al. Missing value estimation methods for DNA microarrays. \textit{Bioinformatics}. 2001;17(6):520--525.
\end{enumerate}

\end{document}

